
\documentclass{article}
\usepackage{colortbl}
\usepackage{makecell}
\usepackage{multirow}
\usepackage{supertabular}

\begin{document}

\newcounter{utterance}

\centering \large Interaction Transcript for game `hot\_air\_balloon', experiment `air\_balloon\_survival\_en\_complexity\_hard', episode 5 with gemini{-}3{-}flash{-}t1.0.
\vspace{24pt}

{ \footnotesize  \setcounter{utterance}{1}
\setlength{\tabcolsep}{0pt}
\begin{supertabular}{c@{$\;$}|p{.15\linewidth}@{}p{.15\linewidth}p{.15\linewidth}p{.15\linewidth}p{.15\linewidth}p{.15\linewidth}}
    \# & $\;$A & \multicolumn{4}{c}{Game Master} & $\;\:$B\\
    \hline

    \theutterance \stepcounter{utterance}  
    & & \multicolumn{4}{p{0.6\linewidth}}{
        \cellcolor[rgb]{0.9,0.9,0.9}{
            \makecell[{{p{\linewidth}}}]{
                \texttt{\tiny{[P1$\langle$GM]}}
                \texttt{You are participating in a collaborative negotiation game.} \\
\\ 
\texttt{Together with another participant, you must agree on a single set of items that will be kept. Each of you has your own view of how much each item matters to you (importance). You do not know how the other participant values the items. Additionally, you are given the effort each item demands.} \\
\texttt{You may only agree on a set if the total effort of the selected items does not exceed a shared limit:} \\
\\ 
\texttt{LIMIT = 7351} \\
\\ 
\texttt{Here are the individual item effort values:} \\
\\ 
\texttt{Item effort = \{"A08": 518, "B95": 656, "A39": 96, "A65": 780, "B81": 190, "B34": 224, "A28": 231, "C82": 181, "B69": 312, "A87": 928, "B64": 826, "B15": 101, "C79": 61, "B16": 814, "C53": 322, "C83": 900, "C10": 744, "A97": 150, "A25": 65, "C84": 853, "A71": 454, "B23": 154, "B96": 237, "C51": 45, "C17": 766, "B19": 294, "C72": 985, "A94": 353, "A46": 60, "B43": 604, "B37": 92, "A88": 454, "A96": 206, "C60": 813, "A75": 234\}} \\
\\ 
\texttt{Here is your personal view on the importance of each item:} \\
\\ 
\texttt{Item importance values = \{"A08": 616, "B95": 754, "A39": 194, "A65": 878, "B81": 288, "B34": 322, "A28": 329, "C82": 279, "B69": 410, "A87": 1026, "B64": 924, "B15": 199, "C79": 159, "B16": 912, "C53": 420, "C83": 998, "C10": 842, "A97": 248, "A25": 163, "C84": 951, "A71": 552, "B23": 252, "B96": 335, "C51": 143, "C17": 864, "B19": 392, "C72": 1083, "A94": 451, "A46": 158, "B43": 702, "B37": 190, "A88": 552, "A96": 304, "C60": 911, "A75": 332\}} \\
\\ 
\texttt{Goal:} \\
\\ 
\texttt{Your goal is to negotiate a shared set of items that benefits you as much as possible (i.e., maximizes total importance to YOU), while staying within the LIMIT. You are not required to make a PROPOSAL in every message {-} you can simply negotiate as well. All tactics are allowed!} \\
\\ 
\texttt{Interaction Protocol:} \\
\\ 
\texttt{You may only use the following structured formats in a message:} \\
\\ 
\texttt{PROPOSAL: \{'A', 'B', 'C', …\}} \\
\texttt{Propose keeping exactly those items.} \\
\texttt{REFUSE: \{'A', 'B', 'C', …\}} \\
\texttt{Explicitly reject opponent's proposal.} \\
\texttt{ARGUMENT: \{'...'\}} \\
\texttt{Defend your last proposal or argue against the player's proposal.} \\
\texttt{AGREE: \{'A', 'B', 'C', …\}} \\
\texttt{Accept the opponent's proposal which ends the game.} \\
\texttt{STRATEGIC REASONING: \{'...'\}} \\
\texttt{	Describe your strategic reasoning or anticipation explaining your choice of action. This is a hidden message which will not be shared with the other participant.} \\
\\ 
\texttt{Rules:} \\
\\ 
\texttt{You may only AGREE on a proposal the other party has logged via PROPOSAL.} \\
\texttt{You may only REFUSE a proposal the other party has logged via PROPOSAL.} \\
\texttt{Total effort of any PROPOSAL or AGREE set must be ≤ LIMIT.} \\
\texttt{Do NOT reveal your hidden importance scores.} \\
\texttt{A tag in a structured format must be followed by colon and whitespace. The argument must be a python set containing 0 or more strings.} \\
\texttt{So, it must be of the form TAG: \{...\}} \\
\texttt{Strictly follow the interaction protocol and DO NOT write anything beyond the given structure.} \\
\texttt{The game ends when one side gives an AGREE to a PROPOSAL made by the other player.} \\
\texttt{The content in your response which can be handed to the other player has to be non{-}empty.} \\
\texttt{Only proposals which have been logged via the PROPOSAL format structure and which haven't been refused via REFUSE are active.} \\
\texttt{You must include the ARGUMENT format at least once somewhere in all of your messages.} \\
\texttt{You must include the STRATEGIC REASONING format only once at the very beginning of every one of your messages and not more often. The contents will not be given to the other player so they can include anything you like including your own importance values. Here you should reason step by step to come up with you next move.} \\
            }
        }
    }
    & \\ \\

    \theutterance \stepcounter{utterance}  
    & \multicolumn{4}{p{0.6\linewidth}}{
        \cellcolor[rgb]{0.9,0.9,0.9}{
            \makecell[{{p{\linewidth}}}]{
                \texttt{\tiny{[P1$\rangle$GM]}}
                \texttt{STRATEGIC REASONING: \{'I will start by identifying the items with the highest value{-}to{-}effort ratio for myself. My top efficiency items are C51, A46, C79, A25, B37, A39, and B15. I will then supplement these with high importance items such as C72, A87, C83, C84, B64, and others until I approach the limit of 7351. I want to establish a strong initial position that includes most of the highly efficient items while staying just under the effort limit. This first proposal will total 7351 effort.'\}} \\
\\ 
\texttt{PROPOSAL: \{'C51', 'A46', 'C79', 'A25', 'B37', 'A39', 'B15', 'A97', 'B23', 'C82', 'B81', 'A96', 'B34', 'A28', 'A75', 'B96', 'B19', 'B69', 'C53', 'A94', 'A71', 'A88', 'A08', 'B43', 'B95', 'C10', 'C17'\}} \\
\\ 
\texttt{ARGUMENT: \{'This selection prioritizes a diverse range of items that balance individual item effort with their overall contribution to the set, while staying exactly within the total effort limit. It includes many of the most efficient choices available.'\}} \\
            }
        }
    }
    & & \\ \\

    \theutterance \stepcounter{utterance}  
    & & \multicolumn{4}{p{0.6\linewidth}}{
        \cellcolor[rgb]{0.9,0.9,0.9}{
            \makecell[{{p{\linewidth}}}]{
                \texttt{\tiny{[GM$\rangle$P2]}}
                \texttt{You are participating in a collaborative negotiation game.} \\
\\ 
\texttt{Together with another participant, you must agree on a single set of items that will be kept. Each of you has your own view of how much each item matters to you (importance). You do not know how the other participant values the items. Additionally, you are given the effort each item demands.} \\
\texttt{You may only agree on a set if the total effort of the selected items does not exceed a shared limit:} \\
\\ 
\texttt{LIMIT = 7351} \\
\\ 
\texttt{Here are the individual item effort values:} \\
\\ 
\texttt{Item effort = \{"A08": 518, "B95": 656, "A39": 96, "A65": 780, "B81": 190, "B34": 224, "A28": 231, "C82": 181, "B69": 312, "A87": 928, "B64": 826, "B15": 101, "C79": 61, "B16": 814, "C53": 322, "C83": 900, "C10": 744, "A97": 150, "A25": 65, "C84": 853, "A71": 454, "B23": 154, "B96": 237, "C51": 45, "C17": 766, "B19": 294, "C72": 985, "A94": 353, "A46": 60, "B43": 604, "B37": 92, "A88": 454, "A96": 206, "C60": 813, "A75": 234\}} \\
\\ 
\texttt{Here is your personal view on the importance of each item:} \\
\\ 
\texttt{Item importance values = \{"A08": 616, "B95": 754, "A39": 194, "A65": 878, "B81": 288, "B34": 322, "A28": 329, "C82": 279, "B69": 410, "A87": 1026, "B64": 924, "B15": 199, "C79": 159, "B16": 912, "C53": 420, "C83": 998, "C10": 842, "A97": 248, "A25": 163, "C84": 951, "A71": 552, "B23": 252, "B96": 335, "C51": 143, "C17": 864, "B19": 392, "C72": 1083, "A94": 451, "A46": 158, "B43": 702, "B37": 190, "A88": 552, "A96": 304, "C60": 911, "A75": 332\}} \\
\\ 
\texttt{Goal:} \\
\\ 
\texttt{Your goal is to negotiate a shared set of items that benefits you as much as possible (i.e., maximizes total importance to YOU), while staying within the LIMIT. You are not required to make a PROPOSAL in every message {-} you can simply negotiate as well. All tactics are allowed!} \\
\\ 
\texttt{Interaction Protocol:} \\
\\ 
\texttt{You may only use the following structured formats in a message:} \\
\\ 
\texttt{PROPOSAL: \{'A', 'B', 'C', …\}} \\
\texttt{Propose keeping exactly those items.} \\
\texttt{REFUSE: \{'A', 'B', 'C', …\}} \\
\texttt{Explicitly reject opponent's proposal.} \\
\texttt{ARGUMENT: \{'...'\}} \\
\texttt{Defend your last proposal or argue against the player's proposal.} \\
\texttt{AGREE: \{'A', 'B', 'C', …\}} \\
\texttt{Accept the opponent's proposal which ends the game.} \\
\texttt{STRATEGIC REASONING: \{'...'\}} \\
\texttt{	Describe your strategic reasoning or anticipation explaining your choice of action. This is a hidden message which will not be shared with the other participant.} \\
\\ 
\texttt{Rules:} \\
\\ 
\texttt{You may only AGREE on a proposal the other party has logged via PROPOSAL.} \\
\texttt{You may only REFUSE a proposal the other party has logged via PROPOSAL.} \\
\texttt{Total effort of any PROPOSAL or AGREE set must be ≤ LIMIT.} \\
\texttt{Do NOT reveal your hidden importance scores.} \\
\texttt{A tag in a structured format must be followed by colon and whitespace. The argument must be a python set containing 0 or more strings.} \\
\texttt{So, it must be of the form TAG: \{...\}} \\
\texttt{Strictly follow the interaction protocol and DO NOT write anything beyond the given structure.} \\
\texttt{The game ends when one side gives an AGREE to a PROPOSAL made by the other player.} \\
\texttt{The content in your response which can be handed to the other player has to be non{-}empty.} \\
\texttt{Only proposals which have been logged via the PROPOSAL format structure and which haven't been refused via REFUSE are active.} \\
\texttt{You must include the ARGUMENT format at least once somewhere in all of your messages.} \\
\texttt{You must include the STRATEGIC REASONING format only once at the very beginning of every one of your messages and not more often. The contents will not be given to the other player so they can include anything you like including your own importance values. Here you should reason step by step to come up with you next move.} \\
\\ 
\texttt{You will now receive the first message of the other player.} \\
\\ 
\texttt{PROPOSAL: \{'C51', 'A46', 'C79', 'A25', 'B37', 'A39', 'B15', 'A97', 'B23', 'C82', 'B81', 'A96', 'B34', 'A28', 'A75', 'B96', 'B19', 'B69', 'C53', 'A94', 'A71', 'A88', 'A08', 'B43', 'B95', 'C10', 'C17'\}} \\
\\ 
\texttt{ARGUMENT: \{'This selection prioritizes a diverse range of items that balance individual item effort with their overall contribution to the set, while staying exactly within the total effort limit. It includes many of the most efficient choices available.'\}} \\
            }
        }
    }
    & \\ \\

    \theutterance \stepcounter{utterance}  
    & & & \multicolumn{4}{p{0.6\linewidth}}{
        \cellcolor[rgb]{0.9,0.9,0.9}{
            \makecell[{{p{\linewidth}}}]{
                \texttt{\tiny{[GM$\langle$P2]}}
                \texttt{STRATEGIC REASONING: \{'The opponent\textbackslash{}'s proposal uses a total effort of 7351 (C51:45 + A46:60 + C79:61 + A25:65 + B37:92 + A39:96 + B15:101 + A97:150 + B23:154 + C82:181 + B81:190 + A96:206 + B34:224 + A28:231 + A75:234 + B96:237 + B19:294 + B69:312 + C53:322 + A94:353 + A71:454 + A88:454 + A08:518 + B43:604 + B95:656 + C10:744 + C17:766 = 7351). My total importance for this set is 8984. I should evaluate if I can do better. High{-}value items missing from their list are C72 (1083), A87 (1026), C83 (998), C84 (951), B64 (924), B16 (912), C60 (911), and A65 (878). These items have high effort but their importance{-}to{-}effort ratio is generally high. For example, C72 (1083/985), A87 (1026/928), C83 (998/900). Let\textbackslash{}'s try to swap some of the lower{-}value high{-}effort items for these. The opponent included C17 (864 importance, 766 effort) and B95 (754 importance, 656 effort). Their list contains a lot of low{-}effort items which is good for filling the limit, but I want to ensure some of the heavy{-}hitters are included. I will propose a set that includes several of my highest value items while maintaining a broad selection to satisfy the opponent.'\}} \\
\\ 
\texttt{REFUSE: \{'C51', 'A46', 'C79', 'A25', 'B37', 'A39', 'B15', 'A97', 'B23', 'C82', 'B81', 'A96', 'B34', 'A28', 'A75', 'B96', 'B19', 'B69', 'C53', 'A94', 'A71', 'A88', 'A08', 'B43', 'B95', 'C10', 'C17'\}} \\
\\ 
\texttt{PROPOSAL: \{'C72', 'A87', 'C83', 'C84', 'B64', 'B16', 'C60', 'A65', 'C17', 'B95', 'C51', 'A46', 'C79', 'A25'\}} \\
\\ 
\texttt{ARGUMENT: \{'I have included many of the most valuable items which also have a high importance{-}to{-}effort ratio. This selection focuses on high{-}impact items while staying within the effort limit. By including these major items, we maximize the total value of our agreement.'\}} \\
            }
        }
    }
     \\ \\

    \theutterance \stepcounter{utterance}  
    & & \multicolumn{4}{p{0.6\linewidth}}{
        \cellcolor[rgb]{0.9,0.9,0.9}{
            \makecell[{{p{\linewidth}}}]{
                \texttt{\tiny{[P1$\langle$GM]}}
                \texttt{REFUSE: \{'C51', 'A46', 'C79', 'A25', 'B37', 'A39', 'B15', 'A97', 'B23', 'C82', 'B81', 'A96', 'B34', 'A28', 'A75', 'B96', 'B19', 'B69', 'C53', 'A94', 'A71', 'A88', 'A08', 'B43', 'B95', 'C10', 'C17'\}} \\
\\ 
\texttt{PROPOSAL: \{'C72', 'A87', 'C83', 'C84', 'B64', 'B16', 'C60', 'A65', 'C17', 'B95', 'C51', 'A46', 'C79', 'A25'\}} \\
\\ 
\texttt{ARGUMENT: \{'I have included many of the most valuable items which also have a high importance{-}to{-}effort ratio. This selection focuses on high{-}impact items while staying within the effort limit. By including these major items, we maximize the total value of our agreement.'\}} \\
            }
        }
    }
    & \\ \\

    \theutterance \stepcounter{utterance}  
    & \multicolumn{4}{p{0.6\linewidth}}{
        \cellcolor[rgb]{0.9,0.9,0.9}{
            \makecell[{{p{\linewidth}}}]{
                \texttt{\tiny{[P1$\rangle$GM]}}
                \texttt{STRATEGIC REASONING: \{'The opponent has proposed a set dominated by high{-}effort, high{-}importance items (C72, A87, C83, C84, B64, B16, C60, A65). My importance values for these are high, but they take up a lot of the effort limit (Total effort of their proposal: 985+928+900+853+826+814+813+780+766+656+45+60+61+65 = 8352). Wait, I must recalculate the effort of their proposal. 985+928+900+853+826+814+813+780+766+656+45+60+61+65 = 8352. The LIMIT is 7351. Their proposal is actually invalid as it exceeds the limit. I will refuse it and propose a set that is within the limit, focusing on many of the items they suggested but trimming the total effort to meet the 7351 constraint while keeping my high{-}value items.'\}} \\
\\ 
\texttt{REFUSE: \{'C72', 'A87', 'C83', 'C84', 'B64', 'B16', 'C60', 'A65', 'C17', 'B95', 'C51', 'A46', 'C79', 'A25'\}} \\
\\ 
\texttt{PROPOSAL: \{'C72', 'A87', 'C83', 'C84', 'B64', 'B16', 'C60', 'C51', 'A46', 'C79', 'A25', 'B37', 'A39', 'B15'\}} \\
\\ 
\texttt{ARGUMENT: \{'Your previous proposal exceeded the effort limit of 7351. I have adjusted the list to include most of the high{-}impact items you prioritized while ensuring we remain within the mandatory effort constraints. This revised list maintains the core of your suggestions but makes it feasible.'\}} \\
            }
        }
    }
    & & \\ \\

    \theutterance \stepcounter{utterance}  
    & & \multicolumn{4}{p{0.6\linewidth}}{
        \cellcolor[rgb]{0.9,0.9,0.9}{
            \makecell[{{p{\linewidth}}}]{
                \texttt{\tiny{[GM$\rangle$P2]}}
                \texttt{REFUSE: \{'C72', 'A87', 'C83', 'C84', 'B64', 'B16', 'C60', 'A65', 'C17', 'B95', 'C51', 'A46', 'C79', 'A25'\}} \\
\\ 
\texttt{PROPOSAL: \{'C72', 'A87', 'C83', 'C84', 'B64', 'B16', 'C60', 'C51', 'A46', 'C79', 'A25', 'B37', 'A39', 'B15'\}} \\
\\ 
\texttt{ARGUMENT: \{'Your previous proposal exceeded the effort limit of 7351. I have adjusted the list to include most of the high{-}impact items you prioritized while ensuring we remain within the mandatory effort constraints. This revised list maintains the core of your suggestions but makes it feasible.'\}} \\
            }
        }
    }
    & \\ \\

    \theutterance \stepcounter{utterance}  
    & & & \multicolumn{4}{p{0.6\linewidth}}{
        \cellcolor[rgb]{0.9,0.9,0.9}{
            \makecell[{{p{\linewidth}}}]{
                \texttt{\tiny{[GM$\langle$P2]}}
                \texttt{STRATEGIC REASONING: \{'The opponent correctly identified that my previous proposal (effort total: 985+928+900+853+826+814+813+780+766+656+45+60+61+65 = 8352) exceeded the limit of 7351. Their counter{-}proposal (C72, A87, C83, C84, B64, B16, C60, C51, A46, C79, A25, B37, A39, B15) has a total effort of 985+928+900+853+826+814+813+45+60+61+65+92+96+101 = 6839. This is well under the limit of 7351 (remaining capacity: 512). My total importance for their set is 1083+1026+998+951+924+912+911+143+158+159+163+190+194+199 = 8011. Since there is 512 capacity left, I should propose adding items to their list to maximize my value. Items I can add within 512 effort: A71 (454 effort, 552 importance), A88 (454 effort, 552 importance), A08 (518 effort {-} too much), A94 (353 effort, 451 importance). Adding A71 seems best.'\}} \\
\\ 
\texttt{REFUSE: \{'C72', 'A87', 'C83', 'C84', 'B64', 'B16', 'C60', 'C51', 'A46', 'C79', 'A25', 'B37', 'A39', 'B15'\}} \\
\\ 
\texttt{PROPOSAL: \{'C72', 'A87', 'C83', 'C84', 'B64', 'B16', 'C60', 'C51', 'A46', 'C79', 'A25', 'B37', 'A39', 'B15', 'A71'\}} \\
\\ 
\texttt{ARGUMENT: \{'I appreciate the correction on the limit. Your last proposal left over 500 units of effort unused. I have added item A71, which utilizes more of our available budget to include another significant item while still remaining comfortably under the 7351 limit.'\}} \\
            }
        }
    }
     \\ \\

    \theutterance \stepcounter{utterance}  
    & & \multicolumn{4}{p{0.6\linewidth}}{
        \cellcolor[rgb]{0.9,0.9,0.9}{
            \makecell[{{p{\linewidth}}}]{
                \texttt{\tiny{[P1$\langle$GM]}}
                \texttt{REFUSE: \{'C72', 'A87', 'C83', 'C84', 'B64', 'B16', 'C60', 'C51', 'A46', 'C79', 'A25', 'B37', 'A39', 'B15'\}} \\
\\ 
\texttt{PROPOSAL: \{'C72', 'A87', 'C83', 'C84', 'B64', 'B16', 'C60', 'C51', 'A46', 'C79', 'A25', 'B37', 'A39', 'B15', 'A71'\}} \\
\\ 
\texttt{ARGUMENT: \{'I appreciate the correction on the limit. Your last proposal left over 500 units of effort unused. I have added item A71, which utilizes more of our available budget to include another significant item while still remaining comfortably under the 7351 limit.'\}} \\
            }
        }
    }
    & \\ \\

    \theutterance \stepcounter{utterance}  
    & \multicolumn{4}{p{0.6\linewidth}}{
        \cellcolor[rgb]{0.9,0.9,0.9}{
            \makecell[{{p{\linewidth}}}]{
                \texttt{\tiny{[P1$\rangle$GM]}}
                \texttt{STRATEGIC REASONING: \{'The opponent correctly pointed out that my previous proposal (effort 6287) left significant unused capacity. Their new proposal adds A71 (effort 454), bringing the total effort to 6741. This is still well under the limit of 7351 (610 units remaining). I want to maximize my importance, so I will add items that have high value{-}to{-}effort for me and fit within the remaining 610 effort. Adding B19 (effort 294) and A96 (effort 206) adds 500 more effort, bringing the total to 7241. This is close to the limit and includes items I value highly.'\}} \\
\\ 
\texttt{REFUSE: \{'C72', 'A87', 'C83', 'C84', 'B64', 'B16', 'C60', 'C51', 'A46', 'C79', 'A25', 'B37', 'A39', 'B15', 'A71'\}} \\
\\ 
\texttt{PROPOSAL: \{'C72', 'A87', 'C83', 'C84', 'B64', 'B16', 'C60', 'C51', 'A46', 'C79', 'A25', 'B37', 'A39', 'B15', 'A71', 'B19', 'A96'\}} \\
\\ 
\texttt{ARGUMENT: \{'We still have room in the budget. By adding B19 and A96, we utilize the effort limit more effectively, bringing our total closer to the maximum allowed while retaining all the high{-}value items we have already agreed upon.'\}} \\
            }
        }
    }
    & & \\ \\

    \theutterance \stepcounter{utterance}  
    & & \multicolumn{4}{p{0.6\linewidth}}{
        \cellcolor[rgb]{0.9,0.9,0.9}{
            \makecell[{{p{\linewidth}}}]{
                \texttt{\tiny{[GM$\rangle$P2]}}
                \texttt{REFUSE: \{'C72', 'A87', 'C83', 'C84', 'B64', 'B16', 'C60', 'C51', 'A46', 'C79', 'A25', 'B37', 'A39', 'B15', 'A71'\}} \\
\\ 
\texttt{PROPOSAL: \{'C72', 'A87', 'C83', 'C84', 'B64', 'B16', 'C60', 'C51', 'A46', 'C79', 'A25', 'B37', 'A39', 'B15', 'A71', 'B19', 'A96'\}} \\
\\ 
\texttt{ARGUMENT: \{'We still have room in the budget. By adding B19 and A96, we utilize the effort limit more effectively, bringing our total closer to the maximum allowed while retaining all the high{-}value items we have already agreed upon.'\}} \\
            }
        }
    }
    & \\ \\

    \theutterance \stepcounter{utterance}  
    & & & \multicolumn{4}{p{0.6\linewidth}}{
        \cellcolor[rgb]{0.9,0.9,0.9}{
            \makecell[{{p{\linewidth}}}]{
                \texttt{\tiny{[GM$\langle$P2]}}
                \texttt{STRATEGIC REASONING: \{'The opponent\textbackslash{}'s new proposal includes C72, A87, C83, C84, B64, B16, C60, C51, A46, C79, A25, B37, A39, B15, A71, B19, and A96. Let\textbackslash{}'s calculate the total effort: 985 + 928 + 900 + 853 + 826 + 814 + 813 + 45 + 60 + 61 + 65 + 92 + 96 + 101 + 454 + 294 + 206 = 7293. This is under the limit of 7351. The remaining effort is 7351 {-} 7293 = 58. My total importance for this set is 8011 (from the previous set) + 552 (A71) + 392 (B19) + 304 (A96) = 9259. There are no items with an effort value of 58 or less remaining except for C51 (45), A46 (60), C79 (61), A25 (65), but those are already in the list. Wait, looking at the effort list, C51 is 45. Is C51 already in? Yes. A46 is 60. Is there anything smaller? No. So 7293 is likely as close as we can get unless we swap items. This set includes most of my top items (C72, A87, C83, C84, B64, B16, C60). I will accept this proposal as it is very strong for me.'\}} \\
\\ 
\texttt{AGREE: \{'C72', 'A87', 'C83', 'C84', 'B64', 'B16', 'C60', 'C51', 'A46', 'C79', 'A25', 'B37', 'A39', 'B15', 'A71', 'B19', 'A96'\}} \\
\\ 
\texttt{ARGUMENT: \{'This proposal effectively utilizes the limit and includes a strong selection of high{-}value items that we both seem to find important. I agree to this set.'\}} \\
            }
        }
    }
     \\ \\

    \theutterance \stepcounter{utterance}  
    & & & \multicolumn{2}{p{0.3\linewidth}}{
        \cellcolor[rgb]{0.9,0.9,0.9}{
            \makecell[{{p{\linewidth}}}]{
                \texttt{\tiny{[GM$|$GM]}}
                \texttt{game successful} \\
            }
        }
    }
    & & \\ \\

    \theutterance \stepcounter{utterance}  
    & & & \multicolumn{2}{p{0.3\linewidth}}{
        \cellcolor[rgb]{0.9,0.9,0.9}{
            \makecell[{{p{\linewidth}}}]{
                \texttt{\tiny{[GM$|$GM]}}
                \texttt{end game} \\
            }
        }
    }
    & & \\ \\

\end{supertabular}
}

\end{document}
